\section{September 9, 2026}

This lecture introduces extrapolation as a way to increase the order of
an existing time-stepping method. We then begin the study of stability
as a property of solutions of the differential equation itself.

\subsection{Accuracy and computational cost}

Consider a method of order \(p\) on the uniform grid
\[
  t_n=t_0+n\tau,
  \qquad n=0,\ldots,N,
  \qquad t_N=T,
\]
and let \(Y_\tau(T)=x_N\) denote the numerical approximation at the
final time. Convergence of order \(p\) means that, for sufficiently
small \(\tau\),
\begin{equation}\label{eq:global-error-sep-09}
  \bigl\lVert Y_\tau(T)-x(T)\bigr\rVert
  \leq C(T)\tau^p.
\end{equation}
In particular, \(Y_\tau(T)\to x(T)\) as \(\tau\to0\). For a fixed step
size, increasing the order generally decreases the error, provided that
\(\tau\) is in the asymptotic regime in which
\eqref{eq:global-error-sep-09} is valid.

Suppose a tolerance \(\varepsilon>0\) is prescribed. Balancing the
error estimate with the tolerance gives
\[
  C(T)\tau^p\approx\varepsilon,
  \qquad
  \tau\approx
  \left(\frac{\varepsilon}{C(T)}\right)^{1/p}.
\]
Since \(N=(T-t_0)/\tau\), the resulting work, measured by the number of
time steps, is approximately
\begin{equation}\label{eq:cost-sep-09}
  N\approx (T-t_0)
  \left(\frac{C(T)}{\varepsilon}\right)^{1/p}.
\end{equation}
Thus a higher-order method has a more favorable dependence on the
tolerance. Formula \eqref{eq:cost-sep-09} does not include the work per
step or the size of the error constant, both of which also matter in
practice.

\subsection{Richardson extrapolation}

Extrapolation requires more than an error bound. Assume the stronger
asymptotic expansion
\begin{equation}\label{eq:error-expansion-sep-09}
  Y_\tau(T)
  =x(T)+c_p(T)\tau^p+c_{p+1}(T)\tau^{p+1}
   +c_{p+2}(T)\tau^{p+2}+\cdots.
\end{equation}
Rerunning the same method with half the step size gives
\[
  Y_{\tau/2}(T)
  =x(T)+c_p(T)\frac{\tau^p}{2^p}
   +c_{p+1}(T)\frac{\tau^{p+1}}{2^{p+1}}
   +O(\tau^{p+2}).
\]
Multiplying this equation by \(2^p\) and subtracting
\eqref{eq:error-expansion-sep-09} cancels the leading error term:
\[
  2^pY_{\tau/2}(T)-Y_\tau(T)
  =(2^p-1)x(T)+O(\tau^{p+1}).
\]
Consequently, the \emph{Richardson extrapolate}
\begin{equation}\label{eq:richardson-sep-09}
  Y^{[1]}_\tau(T)
  =\frac{2^pY_{\tau/2}(T)-Y_\tau(T)}{2^p-1}
\end{equation}
satisfies
\[
  Y^{[1]}_\tau(T)=x(T)+O(\tau^{p+1}).
\]
This gains one order. The underlying one-step method is unchanged; the
additional accuracy comes from combining two complete runs.

\subsubsection{Symmetric methods}

For an autonomous problem, a one-step method is \emph{symmetric} if a
step of length \(-\tau\) exactly reverses a step of length \(\tau\):
\[
  \Psi^{-\tau}=\bigl(\Psi^\tau\bigr)^{-1}.
\]
Symmetric methods have even order. Under the usual smoothness
assumptions, their fixed-time global error expansion contains powers
of one parity. If the method has order \(p\), the expansion takes the
form
\begin{equation}\label{eq:symmetric-expansion-sep-09}
  Y_\tau(T)
  =x(T)+d_p(T)\tau^p+d_{p+2}(T)\tau^{p+2}
   +d_{p+4}(T)\tau^{p+4}+\cdots.
\end{equation}
Applying the same combination as in \eqref{eq:richardson-sep-09}
cancels the \(\tau^p\) term, but now there is no
\(\tau^{p+1}\) term. Hence
\[
  \frac{2^pY_{\tau/2}(T)-Y_\tau(T)}{2^p-1}
  =x(T)+O(\tau^{p+2}),
\]
so one extrapolation raises the order by two.

\subsection{Extrapolation from several step sizes}

The preceding cancellation can be repeated using several numerical
solutions. Choose distinct positive integers \(n_1,\ldots,n_m\), set
\[
  \tau_i=\frac{\bar\tau}{n_i},
  \qquad i=1,\ldots,m,
\]
and compute \(Y_{\tau_i}(T)\) for every \(i\). Motivated by
\eqref{eq:error-expansion-sep-09}, fit the generalized polynomial
\begin{equation}\label{eq:extrapolation-polynomial-sep-09}
  P(h)=a_0+\sum_{j=0}^{m-2}a_{p+j}h^{p+j}
\end{equation}
to these values:
\begin{equation}\label{eq:extrapolation-interpolation-sep-09}
  P(\tau_i)=Y_{\tau_i}(T),
  \qquad i=1,\ldots,m.
\end{equation}
The unknowns \(a_0,a_p,\ldots,a_{p+m-2}\) are determined by the linear
system
\[
  \begin{bmatrix}
    1&\tau_1^p&\tau_1^{p+1}&\cdots&\tau_1^{p+m-2}\\
    1&\tau_2^p&\tau_2^{p+1}&\cdots&\tau_2^{p+m-2}\\
    \vdots&\vdots&\vdots&&\vdots\\
    1&\tau_m^p&\tau_m^{p+1}&\cdots&\tau_m^{p+m-2}
  \end{bmatrix}
  \begin{bmatrix}
    a_0\\a_p\\a_{p+1}\\\vdots\\a_{p+m-2}
  \end{bmatrix}
  =
  \begin{bmatrix}
    Y_{\tau_1}(T)\\Y_{\tau_2}(T)\\\vdots\\Y_{\tau_m}(T)
  \end{bmatrix}.
\]
For a vector-valued ODE, the same coefficient matrix is used for every
component. Extrapolating to zero step size gives
\begin{equation}\label{eq:general-extrapolate-sep-09}
  \overline{x}(T)=P(0)=a_0.
\end{equation}
The first \(m-1\) error terms have been removed, so, under the assumed
full expansion,
\[
  \overline{x}(T)=x(T)+O(\bar\tau^{p+m-1}).
\]

Equivalently, \(a_0\) is a linear combination of the computed values,
\[
  a_0=\sum_{i=1}^m\gamma_iY_{\tau_i}(T),
\]
where the weights are chosen so that
\[
  \sum_{i=1}^m\gamma_i=1,
  \qquad
  \sum_{i=1}^m\gamma_i\tau_i^{p+j}=0,
  \quad j=0,\ldots,m-2.
\]
For \(m=2\), \(n_1=1\), and \(n_2=2\), these conditions recover
\eqref{eq:richardson-sep-09}. For a symmetric error expansion, one
instead uses the powers \(p,p+2,p+4,\ldots\) in
\eqref{eq:extrapolation-polynomial-sep-09}.

\begin{center}
  \begin{tikzpicture}[x=1.35cm,y=1.05cm]
    \draw[->] (0,0) -- (5.3,0) node[right] {$h$};
    \draw[->] (0,0) -- (0,3.1) node[above] {$Y_h(T)$};
    \draw[blue,thick,smooth]
      plot coordinates {(0,1.05) (0.7,1.22) (1.4,1.48)
        (2.2,1.78) (3.0,2.15) (4.0,2.7)};
    \foreach \x/\y/\lab in {1.0/1.34/m,1.8/1.63/3,
                              2.7/2.02/2,3.7/2.53/1} {
      \fill[blue] (\x,\y) circle (1.7pt);
      \draw (\x,0.08) -- (\x,-0.08) node[below] {$\tau_{\lab}$};
    }
    \fill[red] (0,1.05) circle (1.7pt)
      node[left] {$P(0)\approx x(T)$};
    \node[blue,above] at (3.6,2.55) {$P(h)$};
  \end{tikzpicture}
\end{center}

\subsection{Stability of solutions}

Return to the IVP
\begin{equation}\label{eq:ivp-stability-sep-09}
  x'(t)=f(t,x(t)),
  \qquad x(t_0)=x_0,
\end{equation}
and write its solution as \(x(t;t_0,x_0)\) to display its dependence on
the initial data.

\begin{definition}[Stability of a solution]
  The solution through \(x_0\) is \emph{stable} if, for every
  \(\varepsilon>0\), there exists \(\delta>0\) such that
  \[
    \lVert x_0-y_0\rVert<\delta
    \quad\Longrightarrow\quad
    \bigl\lVert x(t;t_0,x_0)-x(t;t_0,y_0)\bigr\rVert<\varepsilon
  \]
  for every \(t\geq t_0\).
\end{definition}

The estimate from Gr\"onwall's inequality in
\eqref{eq:gronwall-stability-aug-26} gives continuous dependence over
every fixed finite interval. Stability is stronger: the same
\(\delta\) must control the separation for all future times.

\begin{definition}[Asymptotic stability]
  A stable solution through \(x_0\) is \emph{asymptotically stable} if
  all sufficiently close initial values \(y_0\) also satisfy
  \[
    \bigl\lVert x(t;t_0,x_0)-x(t;t_0,y_0)\bigr\rVert
    \longrightarrow0
    \qquad\text{as }t\to\infty.
  \]
\end{definition}

For an autonomous equation \(x'=f(x)\), a point \(x_*\) satisfying
\[
  f(x_*)=0
\]
is an \emph{equilibrium}. The solution starting at the equilibrium is
constant: if \(x_0=x_*\), then \(x(t;t_0,x_*)=x_*\) for all
\(t\geq t_0\). Stability of this equilibrium is obtained from the
preceding definitions by comparing this constant solution with
solutions starting at nearby points.

\begin{example}[Linear autonomous system]
  Consider
  \[
    x'=Ax.
  \]
  The origin is an equilibrium, and
  \[
    x(t;t_0,x_0)=e^{(t-t_0)A}x_0.
  \]
  Thus the difference of two solutions is
  \[
    x(t;t_0,x_0)-x(t;t_0,y_0)
    =e^{(t-t_0)A}(x_0-y_0).
  \]
  The system is stable precisely when \(e^{tA}\) remains bounded for
  \(t\geq0\), and it is asymptotically stable precisely when
  \(e^{tA}\to0\). In spectral terms, stability requires every
  eigenvalue to have nonpositive real part and every eigenvalue on the
  imaginary axis to be semisimple. Asymptotic stability holds exactly
  when every eigenvalue has strictly negative real part.

  In the scalar case \(x'=\lambda x\), the solution is
  \(x(t)=e^{\lambda(t-t_0)}x_0\). Hence the equilibrium at the origin
  is stable when \(\operatorname{Re}\lambda\leq0\), and it is
  asymptotically stable when \(\operatorname{Re}\lambda<0\).
\end{example}

These notions describe the behavior of the exact differential
equation. They are distinct from numerical absolute stability, which
asks whether a time-stepping method reproduces the decay of a stable
test equation.
